\section{Discussion and Threats to Validity}
\subsection{What Would Establish Incremental Goal-Graph Value?}
The proposed architecture contains several mechanisms that can improve outcomes for different reasons. A contract can narrow endpoints; provenance can distinguish data sent to the same endpoint; field scope can preserve unrelated metadata; repair can retain the legitimate part of a mixed request. A goal graph may help locate which restrictions are needed, but it should not receive credit for benefits already supplied by those other mechanisms. This is why B2R is a more informative comparator than an unmediated model or a whole-Skill blocker.

In this implementation, the runtime verifies every sensitive action regardless of whether the graph predicted it. The graph's incremental role is to enable repair at selected resources. That role may help avoid unnecessary repair opportunities, but it can also lose utility when reconstruction misses a relevant sink. The experiment therefore tests a particular integration of static selection and dynamic filtering. A lack of advantage would not refute all possible graph-based defenses; it would show that adding this selection layer has not justified its cost under these fixtures and comparators.

The repair optimization itself should be interpreted narrowly. Candidate restrictions are predefined, and their costs favor individual sink restrictions over a global block. Neither the candidate set nor the costs are learned from external human preferences. The minimum is consequently a mathematical property of a small enumerated choice set. It is not evidence that the policy is globally minimal, that the program patch changes the fewest semantics, or that every retained action is necessary. Independent task acceptance supplies the relevant empirical check: whether the restricted execution still achieves the requested result.

\subsection{Internal and Construct Validity}
The task and Skill analyzers have separated inputs, but both use the same model family, and the fixtures were authored with AI assistance. Shared inductive biases may make interpretation easier than it would be for independently collected tasks. Source-span validation protects against unsupported quotations, not against incorrect interpretations of genuine quotations. The contract compiler can still overgrant or undergrant authority. Component measurements therefore accompany, rather than replace, end-to-end outcomes.

The final outcome oracle is stronger than a defense self-report because it inspects complete state and committed events using independently specified fixture expectations. Nevertheless, its expectations come from the same research design as the fixtures. No second human independently labelled task permissions, paths, or success criteria. The code review used a separate AI agent and found important oracle and accounting defects before formal execution; that is software quality assurance, not inter-annotator validation. We cannot report human agreement statistics or treat AI review as a substitute for them.

The closed action language also changes the construct being measured. Operations are explicit, paths are literal, and provenance is attached by trusted tools. Real scripts can hide effects in subprocesses, dependencies, binary assets, indirect callbacks, or remote services. Here, an unsupported operation becomes an error instead of acquiring an alternate operating-system route. This containment makes the experiment safe and interpretable, but it removes entire bypass classes. A claim of unrestricted agent security would exceed the evidence.

False blocking is sensitive to the request stream. If a method blocks a read, subsequent predicted handles may become invalid, and those malformed requests are not independently legal events. If the model never proposes a necessary action, no denial occurs. FBR can therefore look favorable while task success is poor. Reporting both metrics, together with traces and benign completion, is necessary to distinguish a genuinely selective gateway from one that suppresses useful work indirectly.

\subsection{External and Statistical Validity}
The main split holds out named families but shares the base workflow pattern and several mechanisms with development. Two instances vary resource names and contents rather than supplying independent real applications. The three domain labels are contextual descriptions over one substrate. These choices support a controlled mechanism study; they do not establish broad cross-domain generalization. The transfer cohort is also small and changes the original execution environment. Original static bundles provide complementary realism, but classification agreement on text cannot bridge the gap to native execution.

Only \luna{} is tested. Its behavior depends on the API model returned at execution time, the prompt hierarchy, strict output schema, low reasoning effort, and the two-round budget. Other models, longer trajectories, adaptive attackers, or a less explicit instruction hierarchy may produce different attack rates. Repeating calls captures some sampling variation but does not turn one model into a model population. The returned model identifiers and request records are retained so that future replication can identify version differences rather than silently pooling them.

Family-level resampling respects more dependence than treating every method/repeat as independent. Six clusters nevertheless provide limited statistical resolution, and related templates create dependence even across the named clusters. Reported intervals are descriptive uncertainty summaries. A zero attack count under B0 creates an especially important identification limit: there is no observed susceptibility for an added defense to reduce in that cohort. Stress replay can then demonstrate that an enforcement mechanism responds to imposed actions, but cannot manufacture an empirical live-model improvement.

Finally, original benchmark labels encode the benchmark's intended injection conditions, not an independently resolved permission contract for every selected task. Some contextual actions have legitimate uses. Static false positives and false negatives must be interpreted alongside evidence coverage and policy context. We retain incomplete projections and unresolved analyses rather than treating the acquisition process as an invisible preprocessing step.

\subsection{Practical Implications}
A deployable design would need explicit authority bindings, a trusted collector covering every effect, and a way to resolve ambiguous task meaning. The present system's conservative source union prevents easy provenance erasure, but it also loses precision after aggregation. A useful next implementation would preserve field-level or element-level origins through transformations so that authorized parts can be recovered without guessing. That improvement must be tested against both under-tainting attacks and over-tainting utility failures.

Deployment would also require a different response to unresolved authority. This study fails closed automatically because its fixtures have a finite reference specification and no human intervention. In a real workflow, a carefully scoped request for clarification may be preferable to either silently widening access or abandoning the task. Such an interaction would change latency, task success, and user burden, and deserves its own evaluation. We make no claim that the fully automatic protocol establishes the optimal human-facing policy.
